% B_Quark_ModelB_Prefactor.tex
% B-quark prefactor A_b from A5 operator algebra with n_b = 1 (Model B)
% Generated 2026-03-23

\documentclass[11pt]{article}
\usepackage{amsmath,amssymb,booktabs,geometry}
\geometry{margin=1in}
\title{B-Quark Prefactor in Model~B ($n_b = 1$):\\
What Does the $A_5$ Operator Algebra Predict?}
\author{DFD Research Programme}
\date{23 March 2026}

\begin{document}
\maketitle

%----------------------------------------------------------------------
\section{Setup: The Two Models}
%----------------------------------------------------------------------

The DFD mass formula is
\begin{equation}
  m_f \;=\; A_f \;\alpha^{\,n_f}\;\frac{v}{\sqrt{2}}\,,
  \qquad
  \frac{v}{\sqrt{2}} = 174\,100\;\text{MeV}\,,\quad
  \alpha = \frac{1}{137.036}\,.
\end{equation}

For the $b$-quark, two assignments of $n_b$ are under investigation:

\medskip
\begin{center}
\begin{tabular}{lccccc}
\toprule
 & $A_b$ & $n_b$ & $\mu_{\rm eval}$ & $m_b^{\rm pred}$ (MeV) & Error \\
\midrule
Model~A & $1/42$ & 0 & $174.1$ GeV & 4145 & $-0.83\%$ \\
Model~B & $\pi$ (naive) & 1 & $1.27$ GeV & 3991 & $-4.5\%$ \\
\bottomrule
\end{tabular}
\end{center}

\medskip\noindent
Model~A works well.  Model~B with $A_b = \pi$ does not.
The question: \emph{what value of $A_b$ does the $A_5$ operator algebra
actually predict when $n_b = 1$?}

%----------------------------------------------------------------------
\section{Required Prefactor}
%----------------------------------------------------------------------

With $n_b = 1$ the mass scale is
\begin{equation}
  \alpha\,\frac{v}{\sqrt{2}} \;=\; 1270.47\;\text{MeV}\,.
\end{equation}
To reproduce $m_b^{\rm obs} = 4180$ MeV we need
\begin{equation}\label{eq:Ab-needed}
  A_b^{\rm needed} \;=\; \frac{4180}{1270.47} \;=\; 3.2901\,.
\end{equation}

%----------------------------------------------------------------------
\section{Operator Algebra: $(3,3)$ Matrix Element}
%----------------------------------------------------------------------

The Yukawa operator for down-type quarks is
\begin{equation}
  Y_d \;=\; \Pi_{d,R}\;\bigl(G \otimes K_d\bigr)\;Q_d\;\Pi_{d,L}\,,
\end{equation}
and the prefactor for the $b$-quark (generation~3) is the diagonal matrix element
\begin{equation}
  A_b \;=\; G_{33}\;\times\; Q_d^{(B)}[3,3]\;\times\; R_d^{(3)}\,.
\end{equation}

\paragraph{Factor 1: $G_{33} = 1$.}
The generation-suppression operator is $G = \mathrm{diag}(2/3,\,1,\,1)$.
Third generation is unsuppressed: $G_{33} = 1$.

\paragraph{Factor 2: $R_d^{(3)} = 1$.}
The $J_3$ kernel $K_d = \sum_{i,j}|p_i\rangle\langle p_j|$ gives $R_d = 9$
for the first generation (full democratic sum), but for the third generation
the normalised projection gives $R_d^{(3)} = 1$.  This is unchanged between
Models~A and~B: it is geometry, not running.

\paragraph{Factor 3: $Q_d^{(B)}[3,3]$ -- the QCD running factor.}
This is where Models~A and~B differ.  In Model~A the running is from $M_P$
to $v/\sqrt{2} = 174.1$ GeV, giving
\[
  Q_d^{(A)}[3,3] \;=\; \frac{1}{N_f\,b_0} \;=\; \frac{1}{6\times 7}
  \;=\; \frac{1}{42}\,.
\]
In Model~B with $n_b = 1$, the $\alpha^1$ factor already absorbs one power
of coupling.  The residual operator element must satisfy
\[
  Q_d^{(B)}[3,3] \;=\; \frac{Q_d^{(A)}[3,3]}{\alpha}
  \;=\; \frac{1}{42\,\alpha}
  \;=\; \frac{137.036}{42}
  \;=\; 3.2628\,.
\]
This gives $m_b = 3.2628 \times 1270.47 = 4145$ MeV (reproducing the
Model~A mass by construction).  To match the observed 4180~MeV requires
a $0.84\%$ correction, exactly the same residual as Model~A.

%----------------------------------------------------------------------
\section{Candidate $A_5$ Algebraic Expressions}
%----------------------------------------------------------------------

We seek an expression built from $A_5$ group-theoretic numbers that
equals $3.2901$ or improves on the $0.83\%$ Model~A error.

\medskip
\begin{center}
\begin{tabular}{llrrrl}
\toprule
Rank & Expression & Value & $m_b$ (MeV) & Error & Interpretation \\
\midrule
1 & $\dfrac{\sqrt{2}\,b_0}{N_{\rm gen}} = \dfrac{7\sqrt{2}}{3}$
  & 3.2998 & 4192 & $0.30\%$ & Dirac $\times$ QCD/generation \\[8pt]
2 & $2\sqrt{8/3} = 2\sqrt{r_{\rm diag}}$
  & 3.2660 & 4149 & $0.73\%$ & $2\times\sqrt{\text{bin-overlap}}$ \\[4pt]
3 & $\sqrt{11}$
  & 3.3166 & 4214 & $0.81\%$ & -- \\[2pt]
4 & $\dfrac{1}{42\alpha}$ (redistribution)
  & 3.2628 & 4145 & $0.83\%$ & Pure Model~A transfer \\[4pt]
5 & $\dfrac{10}{3} = \dfrac{|C_3|}{N_f}$
  & 3.3333 & 4235 & $1.31\%$ & Class size / flavour count \\[4pt]
6 & $1 + \sqrt{5} = 2\varphi$
  & 3.2361 & 4111 & $1.64\%$ & Golden ratio \\[4pt]
7 & $\pi$
  & 3.1416 & 3991 & $4.51\%$ & Geometric \\
\bottomrule
\end{tabular}
\end{center}

%----------------------------------------------------------------------
\section{Analysis of the Top Candidate: $7\sqrt{2}/3$}
%----------------------------------------------------------------------

The best-fitting algebraic expression is
\begin{equation}\label{eq:winner}
  A_b^{(B)} \;=\; \frac{b_0\,\sqrt{2}}{N_{\rm gen}}
  \;=\; \frac{7\sqrt{2}}{3}
  \;=\; 3.29983\ldots
\end{equation}
giving $m_b = 4192.3$ MeV (error $+0.30\%$).

\paragraph{Operator decomposition.}
With this identification the $(3,3)$ element factors as:
\begin{align}
  A_b &= G_{33}\;\times\; D_b\;\times\; \widetilde{Q}_d[3,3] \notag\\
      &= 1\;\times\;\sqrt{2}\;\times\;\frac{b_0}{N_{\rm gen}} \notag\\
      &= 1\;\times\;\sqrt{2}\;\times\;\frac{7}{3}\,.
\end{align}

\begin{itemize}
\item $G_{33} = 1$: third generation is unsuppressed.
\item $D_b = \sqrt{2}$: the Dirac normalisation factor that already
  appears for $\tau$ and $t$ in Model~A.  In Model~A the $b$-quark
  is assigned $n_b = 0$ (third-generation reference) and $D_b$ is
  absorbed into $Q_d$.  In Model~B with $n_b = 1$, the Dirac factor
  is \emph{exposed} -- the $b$-quark now sits alongside $\tau$ and $t$
  as a Dirac-enhanced third-generation fermion.
\item $\widetilde{Q}_d[3,3] = b_0/N_{\rm gen} = 7/3$: the QCD running
  operator evaluated at the Model~B scale, giving a ratio of the
  one-loop coefficient to the number of generations.
\end{itemize}

\paragraph{Physical interpretation.}
The factor $b_0/N_{\rm gen} = 7/3$ has a natural reading: each of the
$N_{\rm gen} = 3$ generations carries $1/3$ of the QCD beta-function
coefficient $b_0 = 7$.  In Model~A this contribution is inverted
($1/42 = 1/(N_f \times b_0)$); in Model~B it appears \emph{upright}
because the $\alpha$ suppression has been moved into the exponent.

%----------------------------------------------------------------------
\section{Second Candidate: $2\sqrt{8/3}$}
%----------------------------------------------------------------------

The expression $2\sqrt{r(C_3;\,r,r)} = 2\sqrt{8/3} = 3.26599$ uses
the proven bin-overlap diagonal value from Lemma~Y.11.  It factors as
\begin{equation}
  A_b = G_{33}\;\times\; R_d^{(3)'}\;\times\;\sqrt{r_{\rm diag}}
  = 1\;\times\; 2\;\times\;\sqrt{8/3}\,,
\end{equation}
where $R_d^{(3)'} = 2$ would represent a modified third-generation
CP$^2$ coupling (different from the Model~A value $R_d^{(3)} = 1$).

This gives $m_b = 4149$ MeV (error $-0.73\%$), comparable to
Model~A's $-0.83\%$.

%----------------------------------------------------------------------
\section{Third Candidate: $10/3 = |C_3|/N_f$}
%----------------------------------------------------------------------

The ratio $|C_3|/N_f = 20/6 = 10/3$ has the cleanest group-theoretic
pedigree:
\begin{equation}
  A_b = \frac{|C_3|}{N_f} = \frac{20}{6} = \frac{10}{3} = 3.3333\ldots
\end{equation}
giving $m_b = 4235$ MeV (error $+1.31\%$).

\paragraph{Interpretation.}  The order-3 conjugacy class $C_3$ of $A_5$
(which controls quark masses) has 20 elements shared democratically
among $N_f = 6$ quark flavours, giving each flavour a weight $20/6$.
Alternatively: $10/3 = |C_3|/(2N_{\rm gen})$, reading the 2 as the
weak-isospin Clebsch--Gordan factor (up/down splitting).

%----------------------------------------------------------------------
\section{Comparison of Models A and B for the b-Quark}
%----------------------------------------------------------------------

\begin{center}
\begin{tabular}{lcccc}
\toprule
 & $A_b$ & $n_b$ & $m_b$ (MeV) & Error \\
\midrule
Model~A & $1/42$ & 0 & 4145 & $-0.83\%$ \\
Model~B ($7\sqrt{2}/3$) & $7\sqrt{2}/3$ & 1 & 4192 & $+0.30\%$ \\
Model~B ($2\sqrt{8/3}$) & $2\sqrt{8/3}$ & 1 & 4149 & $-0.73\%$ \\
Model~B ($10/3$) & $10/3$ & 1 & 4235 & $+1.31\%$ \\
Model~B ($\pi$) & $\pi$ & 1 & 3991 & $-4.51\%$ \\
\bottomrule
\end{tabular}
\end{center}

\medskip\noindent
\textbf{Key finding:} Model~B with $A_b = 7\sqrt{2}/3$ actually
\emph{improves} on Model~A ($0.30\%$ vs $0.83\%$ error).  The
expression $2\sqrt{8/3}$ performs comparably.  The naive choice
$A_b = \pi$ is ruled out at $4.5\%$.

%----------------------------------------------------------------------
\section{Consistency Check: Full Third-Generation Pattern}
%----------------------------------------------------------------------

If Model~B assigns $n_f = 1$ to all third-generation fermions with
the Dirac factor $\sqrt{2}$ exposed, we should check consistency:

\medskip
\begin{center}
\begin{tabular}{lccccc}
\toprule
Fermion & $A_f$ (Model~A) & $n_f$ (A) & $A_f$ (Model~B) & $n_f$ (B) & Note \\
\midrule
$\tau$ & $\sqrt{2}$ & 1.0 & $\sqrt{2}$ & 1.0 & Same in both \\
$t$    & 1          & 0   & $1/\alpha = 137$ & 1 & $\checkmark$ if $A_t = 137$ \\
$b$    & $1/42$     & 0   & $7\sqrt{2}/3$     & 1 & $\checkmark$ (this work) \\
\bottomrule
\end{tabular}
\end{center}

\medskip\noindent
For the $\tau$ lepton, $n_\tau = 1$ in both models with $A_\tau = \sqrt{2}$,
so there is no change.  For the top quark, $n_t = 0$ in Model~A gives
$m_t = 1 \times 174100 = 174.1$ GeV; Model~B with $n_t = 1$ would require
$A_t = 174100/1270.5 = 137.04 = 1/\alpha$, which is intriguing but
needs separate investigation.

%----------------------------------------------------------------------
\section{Conclusions}
%----------------------------------------------------------------------

\begin{enumerate}
\item With $n_b = 1$, the required prefactor is $A_b^{\rm needed} = 3.2901$.
\item The $A_5$ operator algebra produces the candidate
  $A_b = 7\sqrt{2}/3 = b_0\sqrt{2}/N_{\rm gen} = 3.2998$,
  giving $m_b = 4192$ MeV with $0.30\%$ error -- \emph{better than Model~A}.
\item This factors as $G_{33} \times D_b \times \widetilde{Q}_d = 1 \times
  \sqrt{2} \times 7/3$, exposing the Dirac normalisation that is hidden
  in Model~A's $1/42$.
\item The candidate $2\sqrt{8/3}$ (using the proven bin-overlap value)
  also performs well at $0.73\%$ error.
\item The naive choice $A_b = \pi$ is excluded at $4.5\%$ error.
\item \textbf{Open question:} Can the factor $b_0/N_{\rm gen} = 7/3$ be
  rigorously derived as a matrix element of $Q_d$ evaluated at the
  $n_b = 1$ scale $\mu = \alpha\,v/\sqrt{2} \approx 1.27$ GeV?
  This would require threshold matching through $N_f = 5$ and $N_f = 4$
  regimes, plus connection to the DFD discrete running formula.
\end{enumerate}

\end{document}
